\section{The \pid{} package}
\label{sec:package}
\label{sec:routing}

\pid{} implements \cref{thm:lp} for counterfactual queries under assumptions the user composes. The call takes four objects: finite variables, an observational table, a query, and an assumption set, and returns an interval together with a flag that says whether the interval is the identified set.

For the bow of \cref{ex:bow} under monotone treatment response the call is
\begin{lstlisting}
from partialid import Assumptions, Query, Variables, bound

V = Variables.from_binary("X", "Y")
P = {(0, 0): 0.30, (0, 1): 0.20, (1, 0): 0.15, (1, 1): 0.35}
Q = Query.interventional(outcome="Y", do={"X": 1}, value=1)
a = Assumptions.order("X", "Y") & Assumptions.mtr("Y", "X")
r = bound(V, P, Q, a)
print(r.lower, r.upper, r.tight)      # 0.55 0.85 True
\end{lstlisting}

\subsection{Objects}
\label{sec:design}

\texttt{Variables} declares names and finite domains; variables listed outside \texttt{observed} are latent endogenous variables that the query may still name.

The table is a dictionary of cell probabilities, or the plug-in estimate returned by \texttt{from\_dataframe} or \texttt{from\_counts}.

\texttt{Query} provides interventional probabilities, the ATE, the probabilities of causation (PNS, PN, PS), the counterfactual probability of a single event, and any linear combination of counterfactual events, possibly divided by an observational probability.

\texttt{bound} returns the interval, its width, the flags \texttt{tight} and \texttt{success}, the solver used, and notes on every reduction applied. Three companions share the same signature: \texttt{bound\_path} traces the interval along a sensitivity budget, \texttt{bound\_ci} adds a bootstrap confidence interval for the partially identified parameter \citep{imbens2004confidence}, and \texttt{classify\_program} reports the classification below without solving.

\subsection{Assumption factories}
\label{sec:factories}

Each factory returns a constraint set, and \texttt{\&} intersects them. Zero cells and linear inequalities keep the reduced program linear (\cref{prop:assumptions}); the last three rows of \cref{tab:catalog} change the solver.

\begin{table}[ht]
\centering
\caption{Assumption factories and the form of the constraint each one adds
to the program. Zero cells and linear inequalities keep the reduced program
linear (\cref{prop:assumptions}); the last three rows change the solver.}
\label{tab:catalog}
\small
\begin{tabular}{@{}>{\raggedright\arraybackslash}p{0.30\textwidth}>{\raggedright\arraybackslash}p{0.42\textwidth}l@{}}
\toprule
Factory & Assumption & Constraint \\
\midrule
\texttt{order}, \texttt{diagram}, \texttt{iv} & causal order or diagram
  \citep{manski1990nonparametric,balke1997bounds} & graph \\
\texttt{exclusion} & missing directed edge & zero cells \\
\texttt{mtr}, \texttt{no\_defiers} & monotone response
  \citep{manski1997monotone,imbens1994identification} & zero cells \\
\texttt{mts} & monotone selection \citep{manski2000monotone} & linear inequality \\
\texttt{tan\_msm}, \texttt{rosenbaum\_gamma} & odds-ratio budget, anchored
  \citep{tan2006distributional,rosenbaum2002observational} & linear inequality \\
\texttt{c\_dependence} & propensity budget \citep{masten2018identification} & linear inequality \\
\texttt{margin\_budget} & distance from no confounding & linear inequality \\
\texttt{f\_sensitivity} & $f$-divergence budget \citep{jin2023sensitivity} & convex \\
\texttt{support\_bound} & at most $k$ response types & mixed-integer \\
\texttt{independence} & independent latent blocks & product (outer bound) \\
\bottomrule
\end{tabular}
\end{table}

\subsection{What \texttt{bound} does}
\label{sec:solve}

With the default \texttt{method="auto"}, \texttt{bound} proceeds in three
steps.
\begin{enumerate}[leftmargin=1.4em,itemsep=0.2em]
\item \emph{Point identification.} On an explicit diagram with every
variable observed and no additional shape, sensitivity, or support restrictions, the identification algorithm of
\citet{shpitser2006identification} is run first. If it succeeds, the result
is the identifying functional evaluated on $P$, a degenerate interval.
\item \emph{Classification.} \texttt{classify\_program} restricts the graph
to the iterated closure of the query variables under directed ancestors and
$c$-components (\cref{prop:trim}); pins unconfounded variables
evaluated in a single world (\cref{lem:middle}); drops the $c$-components
the query does not read; and checks the three tests of \cref{thm:lp}.
It also reports the \emph{program degree}, the number of free latent
factors that remain.
\item \emph{Solve.} The class determines the solver and the meaning of the
returned interval (\cref{tab:result}).
\end{enumerate}

\begin{table}[ht]
\centering
\caption{Classification, solver, and what \texttt{bound} returns.
\texttt{tight=False} means that sharpness is not certified, not that the
interval is known to be loose.}
\label{tab:result}
\small
\begin{tabular}{@{}lll@{}}
\toprule
Classification & Solver & Returned interval \\
\midrule
point-identified & ID functional & point, \texttt{tight=True} \\
\cref{thm:lp} holds & closed form or LP (HiGHS) & sharp, \texttt{tight=True} \\
\cref{thm:lp} with a convex budget & convex program (CVXPY) & sharp, \texttt{tight=True} \\
\cref{thm:lp} with a support bound & mixed-integer LP (HiGHS) & sharp, \texttt{tight=True} \\
degree $2$ (bilinear) & nested search, else outer LP & \texttt{tight=False} unless certified \\
degree $\ge3$, \texttt{independence} & LP over merged latents & outer, \texttt{tight=False} \\
assumptions contradict $P$ & --- & \texttt{success=False} \\
\bottomrule
\end{tabular}
\end{table}

Inside the LP path the package merges response types that neither the observational map nor the query distinguishes, and switches to a dual or column-generation formulation when the type space is large. These steps reduce the size of the program without changing its optimum. Closed forms are used for Manski, Balke--Pearl, and a few anchored sensitivity models.

\paragraph{Scope.} The package targets discrete, finite systems. The table $P$ is treated as known; finite-sample uncertainty enters only through \texttt{bound\_ci}. It returns numbers, not symbolic bounds (\cref{sec:related}).

\clearpage
